\section{Geometry-Aware SE(3) Tracking}

\newcommand{\numscales}{3}
\newcommand{\deltadim}{0.22}

\begin{figure}[htbp]
\centering
\resizebox{0.96\textwidth}{!}{%
\begin{tikzpicture}[x=1cm,y=1cm]

  %% ─── Query stack ────────────────────────────────────────────────────────────
  \foreach \i in {\numscales,...,1}{
    \pgfmathsetmacro{\dx}{(\i-1)*\deltadim}
    \pgfmathsetmacro{\dy}{(\i-1)*\deltadim}
    \node[querynode, anchor=north west] (render\i) at ($(\dx,\dy)$) {Render Mesh};
    \node[querynode, anchor=north west] (coords\i) at (render\i.north east) {Coord Buffer};
    \node[querynode, anchor=north west] (feats\i) at (coords\i.north east) {GPU Descriptors};
  }
  \node[font=\scriptsize\sffamily\itshape, text=se3border, right=0.22cm of feats1, align=left] {$L$\\levels};

  \draw[qfeed] ($(render1.west)+(-2.8,0)$) -- (render1.west)
      node[ann, pos=0.45, above] {\scriptsize Query Image/Mesh};

  %% ─── Camera + downscale ─────────────────────────────────────────────────────
  \node[camnode, anchor=north, below=2.0cm of coords1] (DS)
      {Downscale \& Gray\\[-2pt] \scriptsize $(\leq R_{\max}$ px$)$};

  \draw[fw] ($(DS.west)+(-2.6,0)$) -- (DS.west)
      node[ann, pos=0.42, above] {\scriptsize Camera Frame};

  %% ─── State machine ──────────────────────────────────────────────────────────
  \node[smnode, below=0.6cm of DS] (SM)
      {\textbf{State Machine}\\[-2pt] \scriptsize \textsc{matching} $\;\big|\;$ \textsc{optical flow}};
  \draw[fw] (DS.south) -- (SM.north);

  \coordinate (Lcol) at ($(SM.south)+(-3.5,0)$);
  \coordinate (Rcol) at ($(SM.south)+(+3.5,0)$);

  %% ─── LEFT BRANCH — matching ─────────────────────────────────────────────────
  \node[matchnode] (GF) at ($(Lcol)+(0,-1.6)$)
      {GPU Feature Extraction\\[-2pt] \scriptsize Shi-Tomasi $\!\to\!$ orient $\!\to\!$ binary desc};

  \draw[matcharr, rounded corners=4pt]
      (SM.south) -- ++(0,-0.25)
      node[ann, left=4pt, font=\scriptsize\sffamily\itshape, text=se3matchb] {\textsc{matching}}
      -- ++(0,-0.35) -| (GF.north);

  \node[matchnode, below=0.46cm of GF] (MS)
      {Multiscale Match\\[-2pt] \scriptsize $\ell^*\!=\!\arg\max_\ell|\mathcal{M}_\ell|$\enspace Lowe $\tau$};
  \draw[matcharr] (GF.south) -- (MS.north);

  \draw[qfeed, rounded corners=3pt]
      (feats1.south) -- ++(0,-0.5) -- ++(-10.0,0) |- (MS.west);

  \node[matchnode, below=0.46cm of MS] (HG)
      {Homography (RANSAC)\\[-2pt] \scriptsize reprojection $<3$px · quad check};
  \draw[matcharr] (MS.south) -- (HG.north);

  %% ─── RIGHT BRANCH — optical flow ────────────────────────────────────────────
  \node[lknode] (LK) at ($(Rcol)+(0,-1.6)$)
      {Lucas--Kanade Tracking\\[-2pt] \scriptsize prev $\!\to\!$ curr · pyramidal LK};

  \draw[lkarr, rounded corners=4pt]
      (SM.south) -- ++(0,-0.25)
      node[ann, right=4pt, font=\scriptsize\sffamily\itshape, text=se3lkb] {\textsc{optical flow}}
      -- ++(0,-0.35) -| (LK.north);

  \node[lknode, below=0.46cm of LK] (FBV)
      {Fwd--Bwd Validation\\[-2pt] \scriptsize $\|\mathbf{x}_{t-1}\!-\!\hat{\mathbf{x}}_{t-1}\|<\varepsilon$ · boundary};
  \draw[lkarr] (LK.south) -- (FBV.north);

  \node[lknode, below=0.46cm of FBV] (LKHG)
      {Homography (RANSAC)\\[-2pt] \scriptsize LK pts vs.\ query pts · quad check};
  \draw[lkarr] (FBV.south) -- (LKHG.north)
      node[ann, pos=0.5, right=3pt, font=\scriptsize\sffamily] {status pts only};

  %% ─── MERGE ──────────────────────────────────────────────────────────────────
  \coordinate (mergeY) at ($(SM.south |- LKHG.south)+(0,-0.62)$);

  \node[posenode, anchor=north] (LIFT) at (mergeY)
      {2D--3D Lifting\\[-2pt] \scriptsize inliers · coord buffer at $\ell^*$};

  \draw[matcharr, rounded corners=4pt]
      (HG.south) -- ++(0,-0.18) -|
      node[ann, pos=0.65, right=2pt, font=\scriptsize\sffamily] {inliers}
      (LIFT.north west);

  \draw[lkarr, rounded corners=4pt]
      (LKHG.south) -- ++(0,-0.18) -|
      node[ann, pos=0.65, left=2pt, font=\scriptsize\sffamily] {inliers}
      (LIFT.north east);

  \draw[qfeed, rounded corners=3pt]
      (coords1.south) -- ++(0,-0.9) -- ++(-6.0,0) |- (LIFT.west);

  \node[posenode, below=0.48cm of LIFT] (SE3)
      {$\mathrm{SE}(3)$ Gauss--Newton Refinement\\[-2pt] \scriptsize $\mathbf{T}\!\leftarrow\!\exp(\hat{\boldsymbol{\xi}})\mathbf{T}$ \enspace $\min_{\boldsymbol{\xi}}\!\sum_i\|\mathbf{x}_i\!-\!\pi(\mathbf{KTX}_i)\|^2$};
  \draw[pose] (LIFT.south) -- (SE3.north);

  \node[smnode, below=0.48cm of SE3] (TR)
      {Pose Output · confirm\texttt{++}\\[-2pt] \scriptsize $\geq\!\texttt{MIN\_CONFIRM} \Rightarrow$ tracking $=$ \texttt{true} · \textsc{optical flow}};
  \draw[pose] (SE3.south) -- (TR.north);

  %% ─── FAIL NODES (FIXED OVERLAPS) ─────────────────────────────────────────────
  
  %% LEFT BRANCH FAIL (Re-Match)
  % Moved further down and left to clear the Homography box
  \node[warnnode, anchor=north, rotate=90] (FAIL) at ($(HG.south west)+(-2.5,-0.6)$)
      {Re-Match\\[-2pt] \scriptsize quad fail $|$ inliers${<}\theta$ $|$ matches${<}\theta$};

  % Fixed routing: explicitly route around the boxes
  \draw[warn, rounded corners=4pt] 
      (HG.west) -- ++(-0.5,0) |- (FAIL.south)
      node[annW, pos=0.25, above=2pt, rotate=90] {fail};
      
  \draw[warn, rounded corners=4pt] 
      (MS.west) -- ++(-1.2,0) |- (FAIL.south)
      node[annW, pos=0.25, above=2pt, rotate=90] {matches${<}\theta$};
      
  % Return loop: routed completely outside the boxes
  \draw[warn, rounded corners=4pt]
      (FAIL.north) -- ++(-0.6,0) |- (GF.west)
      node[annW, pos=0.25, below=2pt, rotate=90] {tracking $=$ \texttt{false}};


  %% RIGHT BRANCH FAIL (OF Weak / Lost)
  % Moved further down and right to clear the Homography box
  \node[warnnode, anchor=north, rotate=-90] (OFWEAK) at ($(LKHG.south east)+(2.5,-0.6)$)
      {OF Weak / Lost\\[-2pt] \scriptsize pts${<}\texttt{LK\_MIN}$ $|$ quad fail};

  % Fixed routing: explicitly route around the boxes
  \draw[warn, rounded corners=4pt] 
      (LKHG.east) -- ++(0.5,0) |- (OFWEAK.south)
      node[annW, pos=0.25, below=2pt, rotate=-90] {fail};
      
  \draw[warn, rounded corners=4pt] 
      (FBV.east) -- ++(1.2,0) |- (OFWEAK.south)
      node[annW, pos=0.25, below=2pt, rotate=-90] {pts${<}\theta$};
      
  % Return loop: routed completely outside the boxes
  \draw[warn, rounded corners=4pt]
      (OFWEAK.north) -- ++(0.6,0) |- (LK.east)
      node[annW, pos=0.25, below=2pt, rotate=-90] {tracking $=$ \texttt{false} · reset};

  %% ─── BACKGROUND PANELS ──────────────────────────────────────────────────────
  \begin{pgfonlayer}{background}

    %% matching swim-lane
    \node[fill=se3match!30, draw=se3matchb!30, rounded corners=6pt,
          inner sep=12pt, fit=(GF)(MS)(HG)(FAIL)] {};

    %% LK swim-lane
    \node[fill=se3lk!30, draw=se3lkb!30, rounded corners=6pt,
          inner sep=12pt, fit=(LK)(FBV)(LKHG)(OFWEAK)] {};

    %% pose merge lane
    \node[fill=se3orange!25, draw=se3border!25, rounded corners=6pt,
          inner sep=10pt, fit=(LIFT)(SE3)(TR)] {};

    %% outer frame
    \node[fill=gray!6, draw=gray!40, rounded corners=10pt,
          inner sep=22pt,
          fit=(DS)(SM)(GF)(MS)(HG)(LK)(FBV)(LKHG)
              (LIFT)(SE3)(TR)(OFWEAK)(FAIL)] {};

  \end{pgfonlayer}

\end{tikzpicture}%
}
\caption{%
$\mathrm{SE}(3)$ render-and-track pipeline.
\textbf{Top:} A stack of $L$ rendered reference meshes with object-space
coordinate buffers and GPU descriptors.
\textbf{Centre:} The camera frame is downscaled and converted to grayscale,
then routed by the state machine to either the \textcolor{se3matchb}{global
matching branch} (left, blue) or the \textcolor{se3lkb!80!black}{Lucas--Kanade
optical-flow branch} (right, green).
\textbf{Merge:} Inliers from both branches are lifted to 3D via the coordinate
buffer and refined through $\mathrm{SE}(3)$ Gauss--Newton optimisation.
\textbf{Failure handling:} Dashed red paths loop failed matches or lost tracks
back to the corresponding entry points.
}
\label{fig:se3}
\end{figure}

While algebraic tracking based on homography estimation provides efficient frame-to-frame motion estimation, it lacks explicit geometric constraints and becomes unstable under large viewpoint changes, scale variation, or partial occlusions. To address these limitations, we introduce a geometry-aware render-and-track formulation that explicitly estimates pose on the $\mathrm{SE}(3)$ manifold. As illustrated in Figure~\ref{fig:se3}, the system integrates synthetic rendering with a state machine that dynamically switches between global matching and local optical flow.

\subsection{Multiscale Synthetic Reference Rendering}

Instead of relying solely on 2D templates, the proposed method renders a set of $L$ synthetic reference views of the object at multiple depth (scale) hypotheses. For each level, the system generates:
\begin{itemize}
    \item A textured appearance buffer
    \item A descriptor map for feature matching
    \item A dense coordinate buffer storing object-space points
\end{itemize}

Let $\mathbf{X} \in \mathbb{R}^3$ denote a point on the object surface. During rendering, this point is projected into image space using:
\[
\mathbf{x} \sim \mathbf{K} [\mathbf{R} \mid \mathbf{t}] \mathbf{X}
\]
where $\mathbf{K}$ is the intrinsic matrix and $(\mathbf{R}, \mathbf{t}) \in \mathrm{SE}(3)$.

The coordinate buffer encodes $\mathbf{X}$ per pixel, enabling a direct mapping between image coordinates and 3D object geometry.

\subsection{Dual-Branch State Machine}

To ensure efficiency, incoming frames are processed through a two-branch state machine:

\begin{itemize}
    \item \textbf{Global Matching Branch:} Activated during initialization or recovery. Features are extracted and matched against the multiscale reference stack. The optimal scale is selected as:
    \[
    \ell^* = \arg\max_\ell |\mathcal{M}_\ell|
    \]
    where $\mathcal{M}_\ell$ denotes valid matches at level $\ell$.
    
    \item \textbf{Optical Flow Branch:} During stable tracking, correspondences are propagated using pyramidal Lucas--Kanade (LK) optical flow \cite{lucaskanade1981}. Forward-backward consistency is enforced:
    \[
    \|\mathbf{x}_{t-1} - \hat{\mathbf{x}}_{t-1}\| < \varepsilon
    \]
\end{itemize}

In both branches, correspondences are filtered using RANSAC-based homography validation to retain only geometrically consistent inliers.

\subsection{2D--3D Correspondence Lifting}

Each validated feature correspondence is lifted to a 2D--3D pair:
\[
(\mathbf{x}_i, \mathbf{X}_i)
\]
by querying the coordinate buffer. This eliminates the need for triangulation or back-projection and directly provides constraints for pose estimation.

\subsection{Tension-Based Pose Refinement on $\mathrm{SE}(3)$}

Given correspondences $(\mathbf{x}_i, \mathbf{X}_i)$, we define the reprojection function:
\[
\mathbf{p}_i(\mathbf{T}) = \pi(\mathbf{K}\mathbf{T}\mathbf{X}_i)
\]
and residual:
\[
\mathbf{e}_i = \mathbf{x}_i - \mathbf{p}_i(\mathbf{T}), \quad d_i = \|\mathbf{e}_i\|
\]

\paragraph{Spring Model Formulation}

Each correspondence is modeled as a virtual spring in image space:

\begin{itemize}
    \item Current length: $d_i$
    \item Natural length:
    \[
    L = \min_i d_i
    \]
\end{itemize}

The energy of each spring is defined as:
\[
E_i = \frac{\lambda}{2}(d_i - L)^2
\]

\paragraph{Gradient Derivation}

Applying the chain rule:
\[
\frac{dE_i}{d\boldsymbol{\xi}} =
\frac{dE_i}{dd_i}
\cdot
\frac{dd_i}{d\mathbf{e}_i}
\cdot
\frac{d\mathbf{e}_i}{d\boldsymbol{\xi}}
\]

we obtain:
\[
\frac{dE_i}{dd_i} = \lambda(d_i - L), \quad
\frac{dd_i}{d\mathbf{e}_i} = \frac{\mathbf{e}_i^T}{d_i}, \quad
\frac{d\mathbf{e}_i}{d\boldsymbol{\xi}} = -\mathbf{J}_i
\]

Combining:
\[
\nabla E =
\sum_i \mathbf{J}_i^T
\left(
\lambda \frac{d_i - L}{d_i} \mathbf{e}_i
\right)
\]

\paragraph{Weighted Gauss--Newton Formulation}

Defining the weight:
\[
w_i = \frac{d_i - L}{d_i}
\]

the optimization becomes:
\[
E(\mathbf{T}) = \sum_i w_i \|\mathbf{e}_i\|^2
\]

Linearizing:
\[
\mathbf{e}_i(\boldsymbol{\xi} + \delta)
\approx
\mathbf{e}_i + \mathbf{J}_i \delta
\]

leads to the normal equations:
\[
(\mathbf{J}^T \mathbf{W} \mathbf{J}) \delta =
\mathbf{J}^T \mathbf{W} \mathbf{e}
\]

where $\mathbf{W} = \mathrm{diag}(w_i)$.

\paragraph{Pose Update}

The pose is updated using the exponential map:
\[
\mathbf{T} \leftarrow \exp(\delta^\wedge)\mathbf{T}
\]

\paragraph{Interpretation}

This formulation can be interpreted as a physics-inspired weighted least-squares system:
\begin{itemize}
    \item Small residuals ($d_i \approx L$): low influence
    \item Large residuals: strong corrective force
\end{itemize}

Unlike standard Gauss--Newton or fixed robust losses, the weighting is adaptive and derived directly from the residual structure of the current frame.

\subsection{Failure Handling and System Integration}

The system is designed to be fault-tolerant. If RANSAC produces insufficient inliers or the number of tracked LK points falls below a threshold ($\texttt{LK\_MIN}$), tracking is reset and the system re-enters the global matching branch.

Successful pose estimates increment a confirmation counter, and once a minimum threshold ($\geq \texttt{MIN\_CONFIRM}$) is reached, the system resumes efficient optical flow tracking.

\paragraph{Implementation Details}

The rendering stage is implemented using a lightweight cross-platform abstraction built on the \textit{Sokol} framework. On mobile devices, OpenGL ES~3 is used for efficient GPU execution. This ensures that coordinate buffer generation remains portable while fully exploiting available GPU hardware.